\chapter{Allgemeines Vorgehen und Methodik}
\section{Ermittlung der Verbindungen für eine Verbindungsfunktionsstufe $\textrm{VFS}$}
\begin{enumerate}
	\item Initialisierung und Filterung von Bezirken
  \item Bestimmung der Dreiecksverbindungen zwischen den Bezirken durch Delaunay-Triangulation.
  \item Berücksichtigung der $n$-nächsten Nachbarn.
   \item Sicherstellung der Verbindung zu Versorgungszentren.
   \item Filtern der Bezirke und Verbindungen.
   \item Erstellung einer symmetrischen Adjazenzmatrix, die alle Verbindungen korrekt abbildet.
\end{enumerate}

Umsetzung in der Funktion  \mbox{\hyperlink{classluftlinientool_1_1_luftlinien_calculator_a10621c3a62ba5166684e6b769975492e}{calculate\+\_\+vfs}} (self, \mbox{\hyperlink{classluftlinientool_1_1_luftlinien_calculator_a1129397c7e2c387d806eec234eb62c7b}{vfs}}) der Klasse \mbox{\hyperlink{classluftlinientool_1_1_luftlinien_calculator}{Luftlinien\+Calculator}}.

\subsection*{Initialisierung und Filterung von Bezirken}
\begin{itemize}
\item
Zunächst werden die Bezirke mit Stufe $\geq \text{VFS}$ gefiltert, um nur die Bezirke zu berücksichtigen, die für die aktuelle Verbindungsfunktionsstufe $\textrm{VFS}$ relevant sind. Bezirke, die inaktiv sind oder keine Verbindungsanforderungen erfüllen, werden ausgeschlossen.
\item
 Eine Überprüfung erfolgt, um sicherzustellen, dass keine Bezirke identische Koordinaten haben, da dies bei der Delaunay-Triangulation zu Problemen führen könnte.
\end{itemize}

\subsection*{Delaunay-Triangulation}
\begin{itemize}
\item 
Bei genügend aktiven Bezirken wird eine Delaunay-Triangulation durchgeführt, um die Bezirke zu Dreiecken zu verbinden. Diese Dreiecke repräsentieren die Nachbarschaften zwischen den Bezirken. 
\item 
Die resultierenden Nachbarschaftsverbindungen werden in einer Adjazenzmatrix (symmetrisch) gespeichert.
\item 
Kann auch Verbindungen besserer Funktionsstufen ermitteln (Bezirke der Verbindung mit Stufe $>\textrm{VFS}$ ).
\end{itemize}

\subsection*{Berücksichtigung n-nächste Nachbarn}
\begin{itemize}
\item Berechnung Erreichbarkeitsmatrix der Adjazenzmatrix $A$ in $n$ Schritten 
\[
E_n = A^n
\]
\item Update der Adjazenzmatrix = Erreichbarkeitsmatrix.
\end{itemize}


\subsection*{Verbindungen mit Versorgungszentren}
\begin{itemize}
\item
Es wird überprüft, ob jeder Bezirk mit einer ausreichenden Anzahl an Versorgungszentren verbunden ist.
\item
 Wenn dies nicht der Fall ist, werden die nächstgelegenen Versorgungszentren ermittelt und Verbindungen hergestellt.
\end{itemize}

\subsection*{Aktive und inaktive Verbindungen}

\begin{itemize}
\item
Nur aktive Bezirke werden als Quelle und Ziel von Verbindungen berücksichtigt. 
\item
Der Filter „istQuelle“ definiert, welcher Bezirk als Quelle verwendet wird; „istZiel“ definiert die Ziel-Bezirke. Aus Symmetriegründen gilt für die Relationen $od = do$. Dennoch sind „istZiel“-Bezirke kein Teil der Dreiecke, sondern werden von den „istQuelle“-Bezirken erreicht.
\item
Umsetzung: Auf die Adjazenzmatrizen wird eine Filtermaske gelegt. Diese ist das Ergebnis des dyadischen Produkts von „istQuelle“ und „istZiel“ und anschließender Symmetriesierung.

\end{itemize}

Beispiel: Gegeben sei
\glqq istQuelle\grqq{}  $ o = \left[ 
\begin{array}{ccccc}
1 & 1 &0 & 0 & 1 
\end{array} 
\right]^T $ und \glqq istZiel\grqq{} $ d = \left[ \begin{array}{ccccc}
1 & 1 & 1 & 1 & 0 
\end{array} 
\right]^T $. Damit folgt:

\begin{align*}
M_{o \lor d } &= o \cdot d^T &= \left[ 
\begin{array}{ccccc}
1 & 1 & 1 & 1 & 0 \\ 
1 & 1 & 1 & 1 & 0 \\ 
0 & 0 & 0 & 0 & 0 \\ 
0 & 0 & 0 & 0 & 0 \\ 
1 & 1 & 1 & 1 & 0
\end{array} \right] \\
M_{symmetrisch} &= M_{o \lor d } + M_{o \lor d }^T &= 
\left[
\begin{array}{ccccc}
1 & 1 & 1 & 1 & 1 \\ 
1 & 1 & 1 & 1 & 1 \\ 
1 & 1 & 0 & 0 & 1 \\ 
1 & 1 & 0 & 0 & 1 \\ 
1 & 1 & 1 & 1 & 0
\end{array} \right]
\end{align*}

\section{Export der Luftlinieninfrastruktur der VFS}
Der Export der ermittelten Luftlinien der spezifizierten Verbindungsfunktionsstufen als eigenständige Infrastruktur (Knoten, Streckentypen, Strecken, Anbindungen) ermöglicht die Veranschaulichung der Ergebnisse oder die Verwendung des resultierenden Dreiecksnetzes.

Bei der Umwandlung der Adjazenzmatrizen sind zwei Herausforderungen zu beachten:
\begin{itemize}
\item Kompatibilität mit der bereits bestehenden Infrastruktur eines Visummodells.
\item Vermeidung von Überschneidungen bei der Kombination mehrerer Verbindungsfunktionsstufen und korrekte Zuordnung der VFS bei Strecken, die von mehreren Stufen berücksichtigt werden.
\end{itemize}
Deshalb erfolgt der Export der Luftlinieninfrastruktur in 2 Schritten:
\begin{enumerate}
\item Bei Bedarf werden die benötigten Infrastrukturobjekte (Knoten, Streckentypen, Strecken, Anbindungen) unter Berücksichtigung aller berechneten Adjazenzmatrizen aktualisiert und als Attribute der Instanz des Luftlinienkalkulators gespeichert.
Die Umsetzung erfolgt in der Funktion \mbox{\hyperlink{classluftlinientool_1_1_luftlinien_calculator_ae484ccaf33905de92dcc38f3e6f449b7}{extract\+\_\+net}}  der Klasse \mbox{\hyperlink{classluftlinientool_1_1_luftlinien_calculator}{Luftlinien\+Calculator}}.
Die für Visum relevante Nummerierung der Objekte berücksichtigt dabei die vorhandenen Objekte in der aktiven Visuminstanz.
\item Die Funktion \mbox{\hyperlink{classluftlinientool_1_1_luftlinien_calculator_a859aba6233379faafd878b66d51b280c}{export\+\_\+net}} erstellt eine .net Datei mit der Infrastruktur der gewählten Verbindungsfunktionsstufen und öffnet diese in der aktuellen Visuminstanz.
\end{enumerate}

%
%\subsection*{Extrahieren der Infrastruktur}
%
%\subsection*{Export der Infrastruktur }

